\chapter{Introdução} \label{chap:intro}

% Objetivo: apresentar contexto, problema, justificativa, objetivos, hipóteses,
% delimitação e estrutura do trabalho.

\section{Contextualização e motivação}
\textit{Nota terminológica}: as siglas e termos técnicos recorrentes ao longo deste trabalho (TFT, VSN, PICP, MPIW, CRPS, \textit{pinball loss}, \textit{quantile crossing}, rearranjo monotônico, coorte, pré-registro, DM, HLN, MCS, CQR, entre outros) são definidos no Apêndice B (\textit{Glossário técnico de termos recorrentes}). Cada um é introduzido contextualmente em sua primeira ocorrência, com referência cruzada ao glossário quando útil.

A previsão de séries temporais financeiras permanece um problema de elevada relevância acadêmica e aplicada, dada a sua relação direta com decisão de investimento, gestão de risco e alocação de capital. Na literatura econômica, a Hipótese de Mercado Eficiente estabelece um referencial clássico segundo o qual os preços incorporam rapidamente a informação disponível, tornando a previsão sistemática um desafio metodológico relevante \cite{FAMA1991}. Em paralelo, a Hipótese de Mercados Adaptativos propõe que a dinâmica de mercado é evolutiva e dependente de contexto, abrindo espaço para modelos que capturem não linearidades e mudanças de regime \cite{LO2004}.

Uma calibração empírica das expectativas é necessária antes de propor qualquer protocolo de avaliação. Trabalhos de referência em previsão de retornos indicam que o \(R^2\) fora da amostra esperado para retornos de ações líquidas é da ordem de \(0{,}1\%\) a \(1\%\), mesmo com modelos flexíveis e grande conjunto de preditores \cite{GU2020,RAPACH2013}. Esse patamar significa que o componente previsível da média do retorno diário é mínimo: um estudo que aposte seu veredito na acurácia do ponto previsto mede, com alta probabilidade, ruído. A distribuição preditiva, por outro lado, carrega sinal real --- aglomeração de volatilidade, comportamento de caudas e forma da distribuição condicional são mais previsíveis do que seu centro. Por isso, este trabalho adota como objeto científico a \textbf{calibração probabilística} (os quantis previstos correspondem às frequências observadas?) e a \textbf{nitidez} (quão estreitos são os intervalos), e não a acurácia pontual.

Com o aumento de disponibilidade de dados e capacidade computacional, abordagens de aprendizado de máquina passaram a ocupar papel central em previsão financeira, especialmente em cenários multivariados e de múltiplos horizontes \cite{LIMZOHREN2021}. Nesse contexto, o Temporal Fusion Transformer (TFT; ver \hyperref[gl:tft]{Apêndice B}) destaca-se por combinar modelagem temporal, seleção de variáveis, previsão quantílica nativa e interpretabilidade, tornando-se uma referência contemporânea para problemas com covariáveis heterogêneas \cite{LIMETAL2021TFT}.

\section{Problema de pesquisa}
É comum encontrar estudos de previsão financeira avaliados apenas por erro pontual, com o modelo proposto comparado a poucos comparadores simples e com configurações escolhidas após a observação do desempenho fora da amostra. Esse cenário dificulta distinguir ganho real de modelagem de efeito do desenho experimental: a seleção pós-observação infla falsos positivos \cite{WHITE2000,ROMANOWOLF2005}, e o erro pontual é praticamente insensível à qualidade da incerteza estimada.

Sob a ótica de previsão probabilística, avaliar apenas o ponto previsto é insuficiente; é necessário verificar qualidade distributiva, calibração e cobertura de intervalos com regras de pontuação próprias \cite{GNEITING2007}. Previsões quantílicas oferecem um arcabouço consolidado para representar incerteza em diferentes faixas de risco \cite{KOENKERBASSETT1978}, e comparações entre modelos devem respeitar alinhamento temporal e teste estatístico apropriado de acurácia preditiva \cite{DIEBOLDMARIANO1995}. Assim, este trabalho investiga se um TFT quantílico, avaliado sob protocolo pré-registrado e reprodutível, produz previsões calibradas de retornos diários e se supera --- ou empata com --- uma hierarquia de \textit{baselines} explícitos, sem qualquer seleção de configuração por desempenho fora da amostra.

\section{Justificativa}
Este estudo se justifica por três frentes complementares. Na frente científica, desloca a avaliação para o objeto que tem sinal empírico --- a distribuição preditiva --- e integra modelo, dados, métricas probabilísticas e protocolo estatístico em um mesmo desenho confirmatório. Na frente metodológica, adota o TFT como modelo principal por sua adequação a cenários multivariados com múltiplos horizontes e por oferecer mecanismos de interpretação relevantes para análise aplicada \cite{LIMETAL2021TFT}, e o confronta com comparadores fortes (estatísticos e \textit{gradient boosting} quantílico) e com uma referência conformal de cobertura. Na frente prática, a predição quantílica calibrada é insumo direto para decisão sob incerteza e para métricas de risco.

Além disso, ao registrar explicitamente configurações, partições temporais e identificadores de execução, e ao congelar hipóteses e regras de decisão antes da rodada confirmatória, o trabalho se alinha a recomendações contemporâneas de rigor experimental em aprendizado de máquina \cite{PINEAU2021REPRODUCIBILITY,NOSEK2018}.

\section{Objetivos}
\subsection{Objetivo geral}
Caracterizar a calibração probabilística de um Temporal Fusion Transformer aplicado à previsão multi-horizonte de retornos diários, compará-lo a uma hierarquia pré-declarada de \textit{baselines} e descrever a contribuição relativa de famílias de indicadores, sob protocolo estatístico explícito, reprodutível e pré-registrado, no recorte de ativo piloto AAPL e modelos \textit{single-asset}.

\subsection{Objetivos específicos}
\begin{itemize}
    \item Construir um conjunto de dados diário com quatro famílias de \textit{features} (preço, indicadores técnicos, sentimento de notícias e fundamentos), com causalidade temporal verificada por \textit{feature} e mitigação de \textit{leakage}.
    \item Treinar o Temporal Fusion Transformer em modo quantílico sobre uma grade densa de níveis, com rearranjo monotônico dos quantis e rastreabilidade de configurações, sementes e partições.
    \item Definir e aplicar um protocolo experimental com validação \textit{walk-forward} com purga e embargo, conjunto de calibração dedicado, múltiplas sementes e \textit{folds}, alinhamento estrito por \textit{target timestamp} e pré-registro imutável hasheado.
    \item Avaliar a calibração probabilística (cobertura marginal por quantil, PICP, MPIW, \textit{reliability diagram}, cobertura condicional) e o \textit{skill} em \textit{pinball loss} (primária) e CRPS (complementar), comparando o candidato à hierarquia de \textit{baselines} via Diebold-Mariano (HAC/HLN, unilateral) com ajuste de Holm e \textit{Model Confidence Set}, e confrontando a calibração nativa com intervalos conformais (CQR).
    \item Caracterizar a contribuição relativa de famílias de indicadores via \textit{Variable Selection Network}, importância por permutação e ablação, reportando convergências e divergências entre os três métodos.
\end{itemize}

\subsection{Hipóteses}\label{sec:hipoteses}
\begin{itemize}
    \item \textbf{H1 --- Calibração (objeto primário e \textit{gate}).} O candidato produz previsões probabilísticas calibradas de retornos diários --- cobertura marginal por quantil, PICP, \textit{reliability}, nitidez e cobertura condicional --- para ao menos um horizonte, dentro de bandas pré-registradas. H1 é também \textit{gate} de elegibilidade para H2: não se compara \textit{skill} de modelo mal calibrado.
    \item \textbf{H2 --- \textit{Skill} relativo.} Em \textit{pinball loss}, o candidato é comparado a uma hierarquia pré-declarada de \textit{baselines} \{naive, estatísticos, modelo de aprendizado de máquina quantílico\}. Espera-se superar os naive; o interesse científico é superar ou empatar com os fortes; pertencer ao \textit{Model Confidence Set} é evidência complementar; a não-dominância é resultado válido.
    \item \textbf{H3 --- Contribuição de famílias (descritiva).} A contribuição das famílias (preço, técnico, sentimento, fundamento) é heterogênea entre horizontes e detectável de forma consistente em pelo menos dois de três métodos (VSN, permutação, ablação), sem \textit{claim} causal.
\end{itemize}

\section{Delimitação do estudo}
O recorte desta pesquisa é delimitado pela previsão supervisionada de retornos diários com Temporal Fusion Transformer, em desenho \textit{single-asset}: cada modelo é treinado e avaliado para um ativo específico, com AAPL como ativo piloto. O procedimento é replicável para outros ativos, mas a execução confirmatória restringe-se a AAPL. Os horizontes principais de avaliação são \(h+1\) e \(h+7\); \(h+30\) é tratado como suplementar. A avaliação concentra-se em desempenho probabilístico fora da amostra e em comparação pareada contra \textit{baselines} pré-declarados sob alinhamento temporal estrito por \textit{target timestamp}; métricas pontuais são reportadas apenas como perfil descritivo. Toda comparação confirmatória ocorre entre execuções da mesma coorte experimental. Não constitui objetivo deste trabalho comparar exaustivamente arquiteturas profundas, selecionar conjuntos de \textit{features} por desempenho, propor estratégia de \textit{trading}, generalizar para ativos ou períodos não avaliados, nem inferir causalidade entre indicadores e retorno.

\section{Estrutura do trabalho}
Além desta introdução, o Capítulo~\ref{chap:fundamentacao} apresenta a fundamentação teórica sobre previsão financeira, modelagem temporal com aprendizado profundo, predição probabilística e avaliação fora da amostra. O Capítulo~\ref{chap:metodo} descreve o método: dados, engenharia de \textit{features}, configuração do TFT, protocolo experimental, \textit{baselines}, métricas, explicabilidade e regra de decisão. O Capítulo~\ref{chap:resultados} consolida os resultados e a discussão. Por fim, o Capítulo~\ref{chap:conclusoes} sintetiza os achados, explicita limitações e apresenta direções de trabalhos futuros.
